% ============================================================================
\clearpage
\section*{v4.9.0-r45 Source-Only: Chirality Breaks the Reflection of the Triality Form}
\addcontentsline{toc}{section}{v4.9.0-r45 Source-Only: Chirality Breaks the Reflection of the Triality Form}
\label{sec:v490-r45-chirality-breaks-reflection}
% ============================================================================

\paragraph{Purpose.}
Layer v4.9.0-r44 found that mass splitting on the triality slots requires the reflection \(\rho\) to be broken and mixing requires triality
\(\tau\) to be broken.  The author asks whether these breakings follow from the structure.  This layer records the pre-registered test H30
(note \texttt{llm/Vorab\_Notiz\_H30\_Brechung\_aus\_Chiralitaet.md}, commit \texttt{d755a08}, fixed before computing; the identity
\(cL_{e_7}c=-R_{e_7}\) was seen analytically and is disclosed there).  Earlier layers are not rewritten.

\noindent\textbf{Source-only status.}
This v4.9.0-r45 layer is source-only.  It does not build a PDF, does not assign a new DOI,
and does not authorize a public release.  The v4.9.0 Version DOI
\href{https://doi.org/10.5281/zenodo.23534954}{10.5281/zenodo.23534954} and the Concept DOI
\href{https://doi.org/10.5281/zenodo.19210728}{10.5281/zenodo.19210728} are unchanged by it.

\subsection*{v4.9.0-r45 two models}
\label{subsec:v490-r45-models}

Model~I: on \(\mathbb C\otimes\HH\otimes\OO^3\) the chirality projector \(C_{\rm ch}=\tfrac12(1+iL_{e_7})\) of v4.9.0-r42 acts in the same way on
every slot.  Model~II: the chirality operator is \(i(A,B,C)\) with \((A,B,C)\) in the triality algebra \(\mathfrak g\) of v4.9.0-r37 and
\(A=L_{e_7}\); \(B\) and \(C\) are then fixed.  Here \(c\) is octonionic conjugation, and \(\rho(x,y,z)=(\bar x,\bar z,\bar y)\).

\subsection*{v4.9.0-r45 results}
\label{subsec:v490-r45-results}

\paragraph{Theoremlet v4.9.0-r45.1.}
\begin{enumerate}[leftmargin=2em]
\item[(a)] \(cL_{e_7}c=-R_{e_7}\) and \(cR_{e_7}c=-L_{e_7}\).  The plain exchange \((x,y,z)\mapsto(x,z,y)\) is not a symmetry of \(t\); the
transposition symmetries of \(t\) contain \(c\).
\item[(b)] Of the six symmetries \(\tau^k\rho^\varepsilon\) of \(t\), exactly \(1,\tau,\tau^2\) preserve the chirality projector of Model~I; \(\rho\),
\(\tau\rho\), \(\tau^2\rho\) map it to \(\tfrac12(1-iR_{e_7})\), the anomalous choice of v4.9.0-r43.
\item[(c)] \((L_{e_7},L_{e_7},L_{e_7})\notin\mathfrak g\).  The element of \(\mathfrak g\) with first component \(L_{e_7}\) is
\((L_{e_7},R_{e_7},-L_{e_7}-R_{e_7})\); it is not \(\tau\)-invariant.
\item[(d)] In Model~II the chirality operator is \(iR_{e_7}\) on slot~2 and \(-i(L_{e_7}+R_{e_7})\) on slot~3, which is not an involution.
\end{enumerate}

\paragraph{Proof sketch.}
(a) from \(\overline{e_7\bar x}=x\bar e_7=-xe_7\); the rest by explicit matrices and the linear system for \(\mathfrak g\) (script
\texttt{llm/k86\_test\_h30.py}, deviations at most \(10^{-15}\)).  (c): a diagonal element of \(\mathfrak g\) lies in \(\mathfrak g_2\)
(v4.9.0-r37.1(c)), and \(L_{e_7}\) is not a derivation.  \(\square\)

\paragraph{Verdict and reading.}
H30a--d are confirmed.
\begin{enumerate}[leftmargin=2em]
\item The breaking of \(\rho\) follows from the structure.  The transposition symmetries of the triality form contain octonionic
conjugation, conjugation exchanges left and right multiplication, and anomaly cancellation chose the left (v4.9.0-r43).  So \(\rho\) cannot be
a symmetry of the chiral theory; combining it with the frame change \(T\) does not help, since \(T\) changes the sign, not the side.  The
condition of v4.9.0-r44 for mass splitting is therefore met structurally: the slot couplings are generically \(\rho\)-asymmetric, with three
distinct masses in Brannen form.  This also explains why the natural potentials of the earlier mass tests, which did not know chirality, gave
at most two masses.  How strongly \(\rho\) is broken, the phase \(\theta\) and the ratio \(|b|/a\), is not fixed.
\item Triality \(\tau\) survives, but only in Model~I, where the slots are equivalent copies with equal chirality.  The triality-covariant
chirality of Model~II fails: slot~2 would be anomalous and slot~3 would carry no chirality.  The generations therefore cannot be the three
inequivalent triality representations with their own chirality, as read in v4.9.0-r37; they must be equivalent copies, with \(\tau\) permuting
them.  This agrees with the equal quantum numbers of all generations.
\item Mixing stays open: in Model~I \(\tau\) is preserved and, by v4.9.0-r44, there is no mixing.  The breaking needed for mixing has to come
from a different, weaker source, which the structure has not supplied so far.
\end{enumerate}

\subsection*{v4.9.0-r45 status in the three frames}
\label{subsec:v490-r45-frames}

\begin{center}
\small
\begin{tabular}{@{}>{\raggedright\arraybackslash}p{0.30\textwidth}>{\raggedright\arraybackslash}p{0.62\textwidth}@{}}
\toprule
result & M / T / P\\
\midrule
r45.1(a)--(d) & stated in M; in T the chirality is \(-iL_{e_7}\) (v4.9.0-r43.1(d)), and \(\rho\) maps it to \(+iR_{e_7}\), again anomalous: the breaking of \(\rho\) holds in both frames; P not evaluated\\
\bottomrule
\end{tabular}
\end{center}

\subsection*{v4.9.0-r45 register entries}
\label{subsec:v490-r45-register}

\begin{center}
\small
\begin{tabular}{@{}>{\raggedright\arraybackslash}p{0.34\textwidth}>{\raggedright\arraybackslash}p{0.58\textwidth}@{}}
\toprule
statement & status after v4.9.0-r45\\
\midrule
breaking of \(\rho\) & follows from the anomaly-free chirality (r45.1(a), (b)); size not fixed\\
triality \(\tau\) & preserved for equal chirality on all slots (Model~I); a triality-covariant chirality does not exist (r45.1(c), (d))\\
generations as triality slots & only as equivalent copies; the inequivalent representations of v4.9.0-r37 are incompatible with chirality\\
mixing & open; needs a breaking of \(\tau\) from another source\\
\bottomrule
\end{tabular}
\end{center}

\subsection*{v4.9.0-r45 ledger and guard update}
\label{subsec:v490-r45-ledger-guard}

\[
R_{\rm br}^{v490\text{-}r45}=(1,0,0)
\]
for (breaking of \(\rho\) derived from chirality; size of the breaking open; source of the breaking of \(\tau\) open).

The layer spends no guard:
\[
R_{\rm glob}^{(v490\text{-}r45,+)}=(0,1,1,1,1),
\qquad
\Gamma_{\rm guard}^{v490\text{-}r45}=(1,1,1,1,1),
\qquad
\beta_4=1/256\ \text{unspent}.
\]
It does not imply exact finite \(C5\) zero closure, exact all-readout terminality,
\(C6/T2\), a global photon-clock/RH theorem, or global all-readout closure.
